\documentclass[11pt]{article}

\usepackage[margin=1in]{geometry}
\usepackage{amsmath,amssymb,amsthm,mathtools}
\usepackage{booktabs}
\usepackage{enumitem}
\usepackage{hyperref}
\usepackage{microtype}
\usepackage{xcolor}

\hypersetup{colorlinks=true,linkcolor=blue,citecolor=blue,urlcolor=blue}
\allowdisplaybreaks

\newtheorem{proposition}{Proposition}
\newtheorem{lemma}{Lemma}
\newtheorem{definition}{Definition}

\newcommand{\R}{\mathbb R}
\newcommand{\E}{\mathbb E}
\newcommand{\dd}{\,\mathrm d}
\newcommand{\given}{\,\middle|\,}
\newcommand{\ind}{\mathbf 1}
\newcommand{\bfeta}{\boldsymbol\eta}
\newcommand{\bfalpha}{\boldsymbol\alpha}
\newcommand{\bfxi}{\boldsymbol\xi}
\newcommand{\bfz}{\boldsymbol z}
\newcommand{\bfu}{\boldsymbol u}
\newcommand{\bfr}{\boldsymbol r}
\newcommand{\diag}{\operatorname{diag}}
\newcommand{\vecop}{\operatorname{vec}}
\newcommand{\tr}{\operatorname{tr}}

\title{An Annotated Reconstruction of Zhao--Cui Tensor-Train Sequential Filtering}
\author{BayesFilter Technical Note}
\date{2026-06-01}

\begin{document}
\maketitle

\begin{abstract}
This note reconstructs the tensor-train sequential filtering method of Zhao and
Cui in BayesFilter notation.  The goal is not to compress the method, but to
slow it down until the main operations can be followed by a reader who knows
Bayesian state-space models and calculus but has not studied tensor trains.  The
first part follows the paper's state-space model, recursive posterior formula,
basic TT approximation, squared-TT approximation, conditional
Knothe--Rosenblatt maps, particle correction, backward smoothing, and
preconditioning ideas.  The second part isolates a fixed-branch variant needed
for a same-scalar analytical derivative.  Two propositions state precisely what
is proved: the fixed-branch recursion is a normalized approximate filtering
recursion for its own declared approximation, and the fixed-branch derivative
differentiates the same approximate likelihood scalar computed by the forward
pass.
\end{abstract}

\tableofcontents

\section{What Problem Is Being Solved?}

The data arrive in time order.  At time \(t\) there is a hidden state
\[
    x_t\in \R^{m_x}
\]
and an observation
\[
    y_t\in \R^{m_y}.
\]
There may also be a static unknown parameter
\[
    \alpha\in \R^{m_\alpha}.
\]
The filtering problem is to update the conditional density of the current state
and parameter when \(y_t\) arrives:
\[
    p(x_t,\alpha\mid y_{1:t}).
\]
The path problem is to update the density of the whole trajectory and parameter:
\[
    p(x_{0:t},\alpha\mid y_{1:t}).
\]
The smoothing problem asks for an earlier state after later data have been seen:
\[
    p(x_k\mid y_{1:t}),\qquad k<t.
\]

The bottleneck is always an integral.  At time \(t-1\) suppose we know, or have
approximated,
\[
    p(x_{t-1},\alpha\mid y_{1:t-1}).
\]
After seeing \(y_t\), Bayes' rule introduces the unnormalized three-block
density
\[
    q_t(x_t,\alpha,x_{t-1})
    =
    p(x_{t-1},\alpha\mid y_{1:t-1})
    f_\alpha(x_t\mid x_{t-1})
    g_\alpha(y_t\mid x_t).
\]
The next filter is obtained by integrating out \(x_{t-1}\):
\[
    p(x_t,\alpha\mid y_{1:t})
    =
    \frac{\int q_t(x_t,\alpha,x_{t-1})\dd x_{t-1}}
    {\int q_t(x_t,\alpha,x_{t-1})\dd x_t\dd\alpha\dd x_{t-1}}.
\]
In low dimension one can attack this integral by quadrature, particles, or
Gaussian projection.  Zhao--Cui's idea is different: approximate the whole
function \(q_t\) in a separable tensor-train form so that marginal integrals are
cheap contractions instead of high-dimensional numerical integrations.

\section{State-Space Model And BayesFilter Notation}

Zhao--Cui use \(\theta\) for the unknown parameter.  In this note the parameter
is written \(\alpha\), because \(\theta\) is often used later as a generic
argument in derivative formulas.  The paper's model equations (1)--(2) become
\[
    X_t\mid (X_{0:t-1}=x_{0:t-1},\alpha)
    \equiv
    X_t\mid (X_{t-1}=x_{t-1},\alpha)
    \sim f_\alpha(x_t\mid x_{t-1}),
    \tag{BF-1}
\]
and
\[
    Y_t\mid (X_{0:t}=x_{0:t},Y_{1:t-1}=y_{1:t-1},\alpha)
    \equiv
    Y_t\mid (X_t=x_t,\alpha)
    \sim g_\alpha(y_t\mid x_t).
    \tag{BF-2}
\]
The first statement says the latent state process is Markov once \(\alpha\) is
fixed.  The second says the new observation only sees the current state, again
once \(\alpha\) is fixed.

The joint density of parameter, states, and observations through time \(t\),
corresponding to Zhao--Cui Eq. (3), is
\[
    p(\alpha,x_{0:t},y_{1:t})
    =
    p(\alpha)p_\alpha(x_0)
    \prod_{j=1}^{t}
    f_\alpha(x_j\mid x_{j-1})g_\alpha(y_j\mid x_j).
    \tag{BF-3}
\]
This is obtained by the chain rule:
\[
\begin{aligned}
 p(\alpha,x_{0:t},y_{1:t})
 &=
 p(\alpha)p_\alpha(x_0)
 \prod_{j=1}^t
 p(x_j\mid x_{0:j-1},y_{1:j-1},\alpha)
 p(y_j\mid x_{0:j},y_{1:j-1},\alpha)\\
 &=
 p(\alpha)p_\alpha(x_0)
 \prod_{j=1}^t
 f_\alpha(x_j\mid x_{j-1})g_\alpha(y_j\mid x_j),
\end{aligned}
\]
where the last equality uses the Markov and observation assumptions.

Conditioning on data gives Zhao--Cui Eq. (4) in BayesFilter notation:
\[
    p(\alpha,x_{0:t}\mid y_{1:t})
    =
    \frac{p(\alpha,x_{0:t},y_{1:t})}
    {p(y_{1:t})},
    \qquad
    p(y_{1:t})
    =
    \int p(\alpha,x_{0:t},y_{1:t})
    \dd\alpha\,\dd x_{0:t}.
    \tag{BF-4}
\]
The denominator is the evidence.  It is a scalar, but it is usually not
available analytically because the integral spans the parameter and all states.

\section{The Four Marginal Learning Problems}

Zhao--Cui Eqs. (5)--(8) are all marginalizations of the same posterior
\(p(\alpha,x_{0:t}\mid y_{1:t})\).  Filtering keeps only the current state:
\[
    p(x_t\mid y_{1:t})
    =
    \int p(\alpha,x_{0:t}\mid y_{1:t})
    \dd x_{0:t-1}\dd\alpha .
    \tag{BF-5}
\]
Parameter learning keeps only \(\alpha\):
\[
    p(\alpha\mid y_{1:t})
    =
    \int p(\alpha,x_{0:t}\mid y_{1:t})\dd x_{0:t}.
    \tag{BF-6}
\]
Path learning keeps the full state path but removes \(\alpha\):
\[
    p(x_{0:t}\mid y_{1:t})
    =
    \int p(\alpha,x_{0:t}\mid y_{1:t})\dd\alpha .
    \tag{BF-7}
\]
Fixed-lag or retrospective smoothing keeps one old state:
\[
    p(x_k\mid y_{1:t})
    =
    \int p(\alpha,x_{0:t}\mid y_{1:t})
    \dd x_{0:k-1}\dd x_{k+1:t}\dd\alpha ,
    \qquad k<t .
    \tag{BF-8}
\]

The important point is not the names.  It is that the same high-dimensional
posterior density produces all four outputs by integration.  A method that
stores only a cloud of samples can estimate these integrals by sample averages.
A method that stores a Gaussian can compute only Gaussian marginals.  A method
that stores a tensor-train density aims to keep a non-Gaussian function
representation while making many marginal integrals cheap.

\section{The Exact Recursive Bottleneck}

We now derive Zhao--Cui Eqs. (9)--(11).  Start from Bayes' rule:
\[
 p(\alpha,x_{0:t}\mid y_{1:t})
 =
 \frac{p(\alpha,x_{0:t},y_{1:t})}{p(y_{1:t})}.
\]
Using the factorization (BF-3),
\[
 p(\alpha,x_{0:t},y_{1:t})
 =
 p(\alpha,x_{0:t-1},y_{1:t-1})
 f_\alpha(x_t\mid x_{t-1})g_\alpha(y_t\mid x_t).
\]
Also,
\[
 p(\alpha,x_{0:t-1},y_{1:t-1})
 =
 p(\alpha,x_{0:t-1}\mid y_{1:t-1})p(y_{1:t-1}).
\]
Therefore
\[
\begin{aligned}
 p(\alpha,x_{0:t}\mid y_{1:t})
 &=
 \frac{
 p(\alpha,x_{0:t-1}\mid y_{1:t-1})p(y_{1:t-1})
 f_\alpha(x_t\mid x_{t-1})g_\alpha(y_t\mid x_t)}
 {p(y_{1:t})}\\
 &=
 \frac{
 p(\alpha,x_{0:t-1}\mid y_{1:t-1})
 f_\alpha(x_t\mid x_{t-1})g_\alpha(y_t\mid x_t)}
 {p(y_t\mid y_{1:t-1})}.
\end{aligned}
\tag{BF-9}
\]
The last step uses \(p(y_{1:t})=p(y_t\mid y_{1:t-1})p(y_{1:t-1})\).

Now integrate (BF-9) over \(x_{0:t-2}\).  The only factor depending on the old
path \(x_{0:t-2}\) is
\(p(\alpha,x_{0:t-1}\mid y_{1:t-1})\), so
\[
\begin{aligned}
 p(\alpha,x_t,x_{t-1}\mid y_{1:t})
 &=
 \int p(\alpha,x_{0:t}\mid y_{1:t})\dd x_{0:t-2}\\
 &=
 \frac{
 p(\alpha,x_{t-1}\mid y_{1:t-1})
 f_\alpha(x_t\mid x_{t-1})g_\alpha(y_t\mid x_t)}
 {p(y_t\mid y_{1:t-1})}.
\end{aligned}
\tag{BF-10}
\]
Finally, integrate over \(x_{t-1}\) to obtain the next parameter-state filter:
\[
    p(\alpha,x_t\mid y_{1:t})
    =
    \int p(\alpha,x_t,x_{t-1}\mid y_{1:t})\dd x_{t-1}.
    \tag{BF-11}
\]

The numerator in (BF-10) is pointwise evaluable if the previous filter is
evaluable.  The denominator and the marginalization in (BF-11) are the hard
parts.  Zhao--Cui turn this into a function-approximation problem by replacing
the unnormalized numerator with a tensor train.

\section{Tensor Trains From First Principles}

This section starts from the object a numerical reader can picture: a table of
function values.  If \(D=2\), a function \(h(r_1,r_2)\) sampled on a grid is a
matrix.  Low-rank matrix approximation writes
\[
    h(r_1,r_2)\approx \sum_{\rho=1}^{R} u_\rho(r_1)v_\rho(r_2).
\]
Each term is separable: one factor depends only on \(r_1\), one factor only on
\(r_2\).  A tensor train is the many-variable analogue.  It does not require a
full \(D\)-dimensional table.  It stores a chain of low-rank links, so that
information can pass from coordinate \(1\) to coordinate \(D\) through the rank
indices.

Let \(r=(r_1,\ldots,r_D)\in\Omega_1\times\cdots\times\Omega_D\).  A functional
tensor train approximates a scalar function \(h(r)\) by multiplying
matrix-valued one-dimensional functions:
\[
    h(r)
    \approx
    \widehat h(r)
    =
    H_1(r_1)H_2(r_2)\cdots H_D(r_D),
    \tag{TT-1}
\]
where
\[
    H_k(r_k)\in \R^{R_{k-1}\times R_k},\qquad R_0=R_D=1.
\]
Because the first object has one row and the last has one column, the product is
a \(1\times 1\) scalar.  Written with indices,
\[
    \widehat h(r)
    =
    \sum_{\rho_1=1}^{R_1}\cdots
    \sum_{\rho_{D-1}=1}^{R_{D-1}}
    H_1^{1,\rho_1}(r_1)
    H_2^{\rho_1,\rho_2}(r_2)\cdots
    H_D^{\rho_{D-1},1}(r_D).
    \tag{TT-2}
\]
The numbers \(R_k\) are TT ranks.  A high rank lets the representation express
strong interactions across a split \((r_1,\ldots,r_k)\) and
\((r_{k+1},\ldots,r_D)\).  A low rank is useful only when the function has
low-dimensional structure across that split.

To see the role of the ranks, split the variables into a left block and a right
block:
\[
    r_{\le k}=(r_1,\ldots,r_k),\qquad
    r_{>k}=(r_{k+1},\ldots,r_D).
\]
Equation (TT-2) can be regrouped as
\[
    \widehat h(r_{\le k},r_{>k})
    =
    \sum_{\rho=1}^{R_k}
    L_\rho(r_{\le k})R_\rho(r_{>k}),
    \tag{TT-2a}
\]
where
\[
    L_\rho(r_{\le k})
    =
    \sum_{\rho_1,\ldots,\rho_{k-1}}
    H_1^{1,\rho_1}(r_1)\cdots H_k^{\rho_{k-1},\rho}(r_k),
\]
and
\[
    R_\rho(r_{>k})
    =
    \sum_{\rho_{k+1},\ldots,\rho_{D-1}}
    H_{k+1}^{\rho,\rho_{k+1}}(r_{k+1})\cdots
    H_D^{\rho_{D-1},1}(r_D).
\]
Thus \(R_k\) is the number of separated terms allowed across that split.  If a
posterior density has only local or weak long-range dependence in the chosen
ordering, moderate ranks can be enough.  If a posterior has strong global
coupling, the ranks grow and the method becomes expensive.

Each univariate matrix entry is represented in a basis:
\[
    H_k^{a,b}(r_k)
    =
    \sum_{\ell=0}^{p_k-1}
    C_k[a,\ell,b]\psi_{k,\ell}(r_k).
    \tag{TT-3}
\]
The tensor \(C_k\in\R^{R_{k-1}\times p_k\times R_k}\) is the \(k\)-th TT core.
Thus an implementation stores the basis family, the domain map for each
coordinate, and the core arrays \(C_1,\ldots,C_D\).  To evaluate the function at
one point \(r\), it evaluates all basis vectors
\[
    \psi_k(r_k)
    =
    (\psi_{k,0}(r_k),\ldots,\psi_{k,p_k-1}(r_k))^\top,
\]
forms the matrices \(H_k(r_k)\), and multiplies them.

\subsection{A Concrete Basis Choice}

Zhao--Cui's software supports several basis families.  For a self-contained
minimal implementation, choose one.  On a finite interval
\([a_k,b_k]\), map a physical coordinate \(r_k\) to
\[
    z_k=\frac{2r_k-(a_k+b_k)}{b_k-a_k}\in[-1,1],
    \qquad
    r_k=\frac{a_k+b_k}{2}+\frac{b_k-a_k}{2}z_k.
    \tag{TT-6}
\]
Use normalized Legendre polynomials
\[
    \psi_{k,\ell}(r_k)
    =
    \sqrt{\frac{2\ell+1}{b_k-a_k}}\,
    P_\ell(z_k),
    \qquad \ell=0,\ldots,p_k-1.
    \tag{TT-7}
\]
The Legendre polynomials are evaluated by the recurrence
\[
    P_0(z)=1,\qquad P_1(z)=z,
\]
\[
    (\ell+1)P_{\ell+1}(z)
    =
    (2\ell+1)zP_\ell(z)-\ell P_{\ell-1}(z),
    \qquad \ell\ge1.
    \tag{TT-8}
\]
This basis is useful pedagogically because its mass matrix is the identity:
\[
    \int_{a_k}^{b_k}\psi_{k,\ell}(r_k)\psi_{k,\ell'}(r_k)\dd r_k
    =
    \delta_{\ell\ell'}.
    \tag{TT-9}
\]
If another basis is used, every formula below remains the same after replacing
the identity mass matrix by the corresponding one-dimensional mass matrix.

\section{Why TT Marginalization Is Cheap Once The Representation Exists}

Suppose \(h\) is represented by (TT-1).  If we need to integrate out coordinate
\(r_k\), the product structure gives
\[
\begin{aligned}
 \int \widehat h(r_1,\ldots,r_D)\dd r_k
 &=
 H_1(r_1)\cdots H_{k-1}(r_{k-1})
 \left(\int H_k(r_k)\dd r_k\right)
 H_{k+1}(r_{k+1})\cdots H_D(r_D).
\end{aligned}
\tag{TT-4}
\]
Only the \(k\)-th core changes.  Its integrated matrix is
\[
    \overline H_k
    =
    \int H_k(r_k)\dd r_k,\qquad
    \overline H_k[a,b]
    =
    \sum_{\ell=0}^{p_k-1} C_k[a,\ell,b]
    \int \psi_{k,\ell}(r_k)\dd r_k.
    \tag{TT-5}
\]
If the one-dimensional basis integrals are precomputed, marginalizing one
coordinate is a small matrix contraction.  Repeatedly applying (TT-4) integrates
whole blocks of variables.

This explains why Zhao--Cui target a TT representation of
\((x_t,\alpha,x_{t-1})\).  Once that three-block function is separable enough,
the new filter \(p(x_t,\alpha\mid y_{1:t})\) is obtained by integrating the
old-state block \(x_{t-1}\), one coordinate at a time.

\section{Zhao--Cui Algorithm 1, Fully Annotated}

The exact recursion (BF-10) contains the previous exact filter.  A numerical
algorithm has the previous approximation.  Let
\(\widehat\pi_{t-1}(x_{t-1},\alpha)\) denote an unnormalized approximation of
\(p(x_{t-1},\alpha\mid y_{1:t-1})\).  Zhao--Cui Eq. (12) becomes
\[
    q_t(x_t,\alpha,x_{t-1})
    =
    \widehat\pi_{t-1}(x_{t-1},\alpha)
    f_\alpha(x_t\mid x_{t-1})
    g_\alpha(y_t\mid x_t).
    \tag{A1-1}
\]
This is not separable in general.  Even if
\(\widehat\pi_{t-1}\) is a TT over \((\alpha,x_{t-1})\), the transition and
likelihood couple the new and old state coordinates.

Algorithm 1 has three mathematical steps.

\paragraph{Step 1: build a pointwise target.}
Given \(r_t=(x_t,\alpha,x_{t-1})\), evaluate \(q_t(r_t)\) by multiplying the
previous approximation, transition density, and likelihood:
\[
    r_t\mapsto
    \widehat\pi_{t-1}(x_{t-1},\alpha)
    f_\alpha(x_t\mid x_{t-1})
    g_\alpha(y_t\mid x_t).
\]
Storage needed: an evaluator for the previous TT approximation and evaluators
for \(f_\alpha\) and \(g_\alpha\).

\paragraph{Step 2: re-approximate by a TT.}
Choose a variable order, in the paper usually
\[
    (x_t,\alpha,x_{t-1}).
\]
Fit a TT
\[
    \widehat\pi_t(x_t,\alpha,x_{t-1})
    =
    F_{t,1}(x_{t,1})\cdots F_{t,m_x}(x_{t,m_x})
    A_{t,1}(\alpha_1)\cdots A_{t,m_\alpha}(\alpha_{m_\alpha})
    H_{t,1}(x_{t-1,1})\cdots H_{t,m_x}(x_{t-1,m_x})
    \tag{A1-2}
\]
so that \(\widehat\pi_t\approx q_t\).  Zhao--Cui use functional TT tools with
alternating approximation and rank adaptation.  For implementation one must
choose a domain map, basis, fitting points, rank rule, and stopping rule.  The
paper's research code supplies one such family; the fixed-branch section below
supplies a differentiable BayesFilter specialization.

The most direct implementation view is least squares.  Choose fitting points
\[
    r^{(j)}=(x_t^{(j)},\alpha^{(j)},x_{t-1}^{(j)}),
    \qquad j=1,\ldots,N_{\rm fit},
\]
and define target values
\[
    y_j=q_t(r^{(j)}).
\]
For fixed TT ranks and fixed all cores except core \(k\), the fitted function is
linear in the entries of core \(C_k\).  There is a design row \(A_{j,k}\) such
that
\[
    \widehat\pi_t(r^{(j)})=A_{j,k}g_k,
    \qquad g_k=\vecop(C_k).
    \tag{A1-2a}
\]
Let \(L_{j,k}\) be the left environment obtained by multiplying all cores before
\(k\) at point \(r^{(j)}\), and let \(R_{j,k}\) be the right environment from
all cores after \(k\).  Let
\(\psi_k(r_k^{(j)})\) be the local basis vector.  Then
\[
    A_{j,k}
    =
    R_{j,k}^{\top}\otimes \psi_k(r_k^{(j)})^\top\otimes L_{j,k}.
    \tag{A1-2b}
\]
This formula is worth lingering over.  The left environment tells the current
core what the earlier coordinates have already contributed.  The right
environment tells it what later coordinates will contribute.  The basis vector
gives the local coordinate dependence.  The Kronecker product stitches these
three pieces into a row that multiplies the vectorized core.

With weights \(w_j>0\), the core update solves
\[
    \min_{g_k}
    \sum_{j=1}^{N_{\rm fit}}w_j\bigl(A_{j,k}g_k-y_j\bigr)^2
    +\rho\|g_k\|_2^2,
    \tag{A1-2c}
\]
or equivalently
\[
    (A_k^\top W A_k+\rho I)g_k=A_k^\top Wy.
    \tag{A1-2d}
\]
An alternating scheme sweeps through the cores, updating one core at a time.
Zhao--Cui's adaptive implementation is more sophisticated, but this least-
squares view is the easiest way to see how a TT approximation becomes an
actual numerical object.

\paragraph{Step 3: integrate the old state.}
The next approximate filter is
\[
    \widehat\pi_t(x_t,\alpha)
    =
    \int \widehat\pi_t(x_t,\alpha,x_{t-1})\dd x_{t-1}.
    \tag{A1-3}
\]
Because \(x_{t-1}\) occupies the last block, this integral is a sequence of core
integrals:
\[
    \widehat\pi_t(x_t,\alpha)
    =
    F_{t,1}(x_{t,1})\cdots F_{t,m_x}(x_{t,m_x})
    A_{t,1}(\alpha_1)\cdots A_{t,m_\alpha}(\alpha_{m_\alpha})
    \overline H_{t,1}\cdots \overline H_{t,m_x}.
    \tag{A1-4}
\]
The approximate normalizer is
\[
    \widehat Z_t
    =
    \int \widehat\pi_t(x_t,\alpha,x_{t-1})
    \dd x_t\,\dd\alpha\,\dd x_{t-1},
    \tag{A1-5}
\]
again a complete TT contraction.

The main failure mode of Algorithm 1 is sign.  A fitted TT is a real-valued
function approximation.  Truncation or interpolation can make
\(\widehat\pi_t(r_t)<0\), which is impossible for a density.  Negative values
also break transport maps and importance-sampling ratios.  This motivates the
squared-TT construction.

\section{Why Nonnegativity Fails And Why Square-Root TT Repairs It}

Approximating a nonnegative function by a truncated expansion does not preserve
nonnegativity.  The one-dimensional analogy is simple:
\[
    h(x)\ge 0
    \quad\text{but}\quad
    \widehat h(x)=\sum_{\ell=0}^{p-1}c_\ell\psi_\ell(x)
    \text{ may be negative.}
\]
The tensor-train case has the same issue.  A low-rank approximation can reduce
error in a least-squares sense while still creating small negative regions.

Zhao--Cui's repair is to approximate the square root of an unnormalized density.
If
\[
    \sqrt{\pi(r)}\approx \phi(r),
\]
then
\[
    \phi(r)^2\ge 0
\]
for all \(r\).  A small defensive density is then added so that the proposal
density is positive wherever the target might be positive.

\section{Squared-TT Density, Defensive Reference, And Normalizer}

Zhao--Cui Eq. (13) becomes the following.  Let \(\pi(r)\ge0\) be an
unnormalized target on a domain \(\Omega\subseteq\R^D\), with normalizer
\[
    Z=\int_\Omega \pi(r)\dd r.
\]
Let \(\lambda(r)\) be a normalized reference density on \(\Omega\),
\[
    \lambda(r)\ge0,\qquad \int_\Omega \lambda(r)\dd r=1,
\]
and choose \(\tau>0\).  Fit a TT \(\phi(r)\) to \(\sqrt{\pi(r)}\).  Define
\[
    \widehat\pi(r)
    =
    \phi(r)^2+\tau\lambda(r),
    \qquad
    \widehat Z
    =
    \int_\Omega \widehat\pi(r)\dd r,
    \qquad
    \widehat p(r)
    =
    \frac{\widehat\pi(r)}{\widehat Z}.
    \tag{S1}
\]
Since \(\lambda\) is normalized,
\[
    \widehat Z
    =
    \int_\Omega \phi(r)^2\dd r+\tau.
    \tag{S2}
\]
The defensive term has two jobs.  First, it keeps the approximation positive
even if \(\phi\) is zero somewhere.  Second, if \(\pi/\lambda\) is bounded, then
\[
    \frac{p(r)}{\widehat p(r)}
    =
    \frac{\pi(r)/Z}{\widehat\pi(r)/\widehat Z}
    =
    \frac{\widehat Z}{Z}
    \frac{\pi(r)}{\phi(r)^2+\tau\lambda(r)}
    \le
    \frac{\widehat Z}{Z\tau}
    \frac{\pi(r)}{\lambda(r)}.
    \tag{S3}
\]
This is the importance-sampling reason for the defensive density: the true
target is absolutely continuous with respect to the approximation when the
reference covers the target support.

\section{Squared-TT Marginalization And Mass Matrices}

The square changes the marginalization algebra.  Let
\[
    \phi(r)=H_1(r_1)\cdots H_D(r_D),
    \qquad H_k(r_k)\in\R^{R_{k-1}\times R_k}.
\]
For a split after coordinate \(k\), define
\[
    H_{\le k}(r_{\le k})=H_1(r_1)\cdots H_k(r_k)\in\R^{1\times R_k},
\]
and
\[
    H_{>k}(r_{>k})=H_{k+1}(r_{k+1})\cdots H_D(r_D)\in\R^{R_k\times 1}.
\]
Then
\[
    \phi(r)=H_{\le k}(r_{\le k})H_{>k}(r_{>k}).
\]
The marginal square term is
\[
\begin{aligned}
 \int \phi(r)^2\dd r_{>k}
 &=
 \int
 \left[
 H_{\le k}(r_{\le k})H_{>k}(r_{>k})
 \right]^2\dd r_{>k}\\
 &=
 \sum_{a=1}^{R_k}\sum_{b=1}^{R_k}
 H_{\le k,a}(r_{\le k})
 \left[
 \int H_{>k,a}(r_{>k})H_{>k,b}(r_{>k})\dd r_{>k}
 \right]
 H_{\le k,b}(r_{\le k}).
\end{aligned}
\tag{M1}
\]
Define the right mass matrix
\[
    M_{>k}[a,b]
    =
    \int H_{>k,a}(r_{>k})H_{>k,b}(r_{>k})\dd r_{>k}.
    \tag{M2}
\]
This matrix is positive semidefinite.  If \(M_{>k}=L_{>k}L_{>k}^\top\) is a
Cholesky or square-root factorization, then
\[
    \int \phi(r)^2\dd r_{>k}
    =
    \sum_{\gamma=1}^{R_k}
    \left[
    \sum_{a=1}^{R_k}H_{\le k,a}(r_{\le k})L_{>k}[a,\gamma]
    \right]^2.
    \tag{M3}
\]
Including the defensive reference, if
\(\lambda(r)=\prod_{j=1}^D\lambda_j(r_j)\), gives
\[
    \widehat p(r_{\le k})
    =
    \frac{
    \sum_{\gamma=1}^{R_k}
    \left[
    \sum_{a=1}^{R_k}H_{\le k,a}(r_{\le k})L_{>k}[a,\gamma]
    \right]^2
    +
    \tau\prod_{j=1}^k\lambda_j(r_j)}
    {\widehat Z}.
    \tag{M4}
\]
This is Zhao--Cui Eq. (14) in BayesFilter notation.

The practical recursion starts from the right end.  Set \(M_{>D}=1\).  Moving
from \(j=D\) down to \(1\), compute
\[
    M_{>j-1}
    =
    \int H_j(r_j)M_{>j}H_j(r_j)^\top\dd r_j.
    \tag{M5}
\]
In basis coefficients this integral is finite-dimensional.  If
\[
    H_j^{a,b}(r_j)=\sum_{\ell}C_j[a,\ell,b]\psi_{j,\ell}(r_j),
\]
and
\[
    B_j[\ell,\ell']
    =
    \int \psi_{j,\ell}(r_j)\psi_{j,\ell'}(r_j)\dd r_j,
\]
then
\[
    M_{>j-1}[a,a']
    =
    \sum_{b,b'=1}^{R_j}
    \sum_{\ell,\ell'}
    C_j[a,\ell,b]\,M_{>j}[b,b']\,C_j[a',\ell',b']\,B_j[\ell,\ell'].
    \tag{M6}
\]
This formula is the bridge from the paper to code: marginalization stores mass
matrices or their square roots and never enumerates a full tensor grid.

\subsection{A Small Three-Coordinate Marginalization}

Take \(D=3\), and suppose
\[
    \phi(r_1,r_2,r_3)=H_1(r_1)H_2(r_2)H_3(r_3).
\]
If \(H_1\in\R^{1\times R_1}\),
\(H_2\in\R^{R_1\times R_2}\), and
\(H_3\in\R^{R_2\times 1}\), then marginalizing out \(r_3\) gives
\[
\begin{aligned}
 \int \phi(r_1,r_2,r_3)^2\dd r_3
 &=
 H_1(r_1)H_2(r_2)
 \left[\int H_3(r_3)H_3(r_3)^\top\dd r_3\right]
 H_2(r_2)^\top H_1(r_1)^\top.
\end{aligned}
\tag{M7}
\]
The bracketed quantity is an \(R_2\times R_2\) matrix.  It is not a function of
\(r_1\) or \(r_2\); it is a stored contraction.  If we also marginalize out
\(r_2\), then
\[
    M_{>1}
    =
    \int
    H_2(r_2)
    \left[\int H_3(r_3)H_3(r_3)^\top\dd r_3\right]
    H_2(r_2)^\top
    \dd r_2,
    \tag{M8}
\]
and the one-coordinate marginal is
\[
    \int\phi(r_1,r_2,r_3)^2\dd r_2\dd r_3
    =
    H_1(r_1)M_{>1}H_1(r_1)^\top.
    \tag{M9}
\]
This is exactly what the recursive mass formula (M5) does.  The notation in the
paper is compact because it assumes the reader already sees this matrix chain.
The implementation sees it as repeated multiplication and integration of small
matrices.

\section{Conditional Densities And KR Maps From Marginal Ratios}

For a normalized density \(\widehat p(r_1,\ldots,r_D)\), the chain rule gives
\[
    \widehat p(r)
    =
    \widehat p(r_1)
    \prod_{k=2}^D
    \widehat p(r_k\mid r_{<k}),
    \qquad
    \widehat p(r_k\mid r_{<k})
    =
    \frac{\widehat p(r_{\le k})}{\widehat p(r_{<k})}.
    \tag{K1}
\]
The squared-TT marginal formula (M4) supplies both numerator and denominator.
Thus each conditional density is computable from TT contractions.

Define conditional distribution functions
\[
    F_1(r_1)=\int_{-\infty}^{r_1}\widehat p(s_1)\dd s_1,
    \tag{K2}
\]
and, for \(k\ge2\),
\[
    F_k(r_k\mid r_{<k})
    =
    \int_{-\infty}^{r_k}
    \widehat p(s_k\mid r_{<k})\dd s_k.
    \tag{K3}
\]
The lower triangular Knothe--Rosenblatt map is
\[
    F(r)
    =
    \bigl(
    F_1(r_1),
    F_2(r_2\mid r_1),
    \ldots,
    F_D(r_D\mid r_{<D})
    \bigr).
    \tag{K4}
\]
If \(R\sim\widehat p\), then \(F(R)\) is uniform on \([0,1]^D\).  Conversely, if
\(\Xi\sim \mathrm{Unif}([0,1]^D)\), then \(F^{-1}(\Xi)\sim\widehat p\).

The proof is the triangular Jacobian calculation.  The Jacobian of \(F\) is
lower triangular, and
\[
    \frac{\partial F_k(r_k\mid r_{<k})}{\partial r_k}
    =
    \widehat p(r_k\mid r_{<k}).
\]
Therefore
\[
    |\det\nabla F(r)|
    =
    \prod_{k=1}^D \widehat p(r_k\mid r_{<k})
    =
    \widehat p(r).
    \tag{K5}
\]
The pullback of the uniform density by \(F\) is exactly
\(|\det\nabla F(r)|=\widehat p(r)\).

\section{Zhao--Cui Algorithm 2, Fully Annotated}

Algorithm 2 is Algorithm 1 with a square-root TT and a defensive density.
At time \(t-1\), assume the previous filter approximation is evaluable as
\[
    \widehat p_{t-1}(x_{t-1},\alpha).
\]
The pointwise unnormalized target, Zhao--Cui Eq. (15), is
\[
    q_t(x_t,\alpha,x_{t-1})
    =
    \widehat p_{t-1}(x_{t-1},\alpha)
    f_\alpha(x_t\mid x_{t-1})
    g_\alpha(y_t\mid x_t).
    \tag{A2-1}
\]
Fit a TT to the square root:
\[
    \sqrt{q_t(x_t,\alpha,x_{t-1})}
    \approx
    \phi_t(x_t,\alpha,x_{t-1}).
    \tag{A2-2}
\]
Then define Zhao--Cui Eq. (16):
\[
    \widehat q_t(x_t,\alpha,x_{t-1})
    =
    \phi_t(x_t,\alpha,x_{t-1})^2
    +
    \tau_t\lambda_t(x_t)\lambda_\alpha(\alpha)\lambda_{t-1}(x_{t-1}).
    \tag{A2-3}
\]
The approximate normalizer is
\[
    \widehat Z_t
    =
    \int \widehat q_t(x_t,\alpha,x_{t-1})
    \dd x_t\,\dd\alpha\,\dd x_{t-1}.
    \tag{A2-4}
\]
The next normalized approximate filter is
\[
    \widehat p_t(x_t,\alpha)
    =
    \frac{\int \widehat q_t(x_t,\alpha,x_{t-1})\dd x_{t-1}}
    {\widehat Z_t}.
    \tag{A2-5}
\]

Everything is now nonnegative.  The approximation risk has moved: the method
must fit the square root accurately enough, choose a defensive term that is not
too small for stability or too large for accuracy, and keep TT ranks manageable.

\section{Forward Conditional Map And Particle-Filter Correction}

The same \(\widehat q_t(x_t,\alpha,x_{t-1})\) defines a conditional density for
the new state given the old state and parameter:
\[
    \widehat p_t(x_t\mid \alpha,x_{t-1},y_{1:t})
    =
    \frac{\widehat q_t(x_t,\alpha,x_{t-1})}
    {\int \widehat q_t(s,\alpha,x_{t-1})\dd s}.
    \tag{F1}
\]
Constructing the corresponding upper triangular conditional KR map gives a
proposal sampler:
\[
    \widehat X_t^{(i)}
    =
    (F_{t,t}^{u})^{-1}
    \left(\Xi_t^{(i)}\mid \widehat\alpha^{(i)},\widehat x_{t-1}^{(i)}\right),
    \qquad
    \Xi_t^{(i)}\sim\mathrm{Unif}([0,1]^{m_x}).
    \tag{F2}
\]
This is Zhao--Cui Eq. (21) in notation adapted to BayesFilter.

The proposal density for the full path after sampling is
\[
    \widehat p_t(x_t\mid \alpha,x_{t-1},y_{1:t})
    p(\alpha,x_{0:t-1}\mid y_{1:t-1}),
    \tag{F3}
\]
which corresponds to Zhao--Cui Eq. (22).  The exact posterior recursion is
\[
    p(\alpha,x_{0:t}\mid y_{1:t})
    \propto
    p(\alpha,x_{0:t-1}\mid y_{1:t-1})
    f_\alpha(x_t\mid x_{t-1})
    g_\alpha(y_t\mid x_t).
\]
Dividing exact target by proposal gives the unnormalized correction factor
\[
    \omega_f(\alpha,x_{t-1:t})
    =
    \frac{
    f_\alpha(x_t\mid x_{t-1})g_\alpha(y_t\mid x_t)}
    {\widehat p_t(x_t\mid \alpha,x_{t-1},y_{1:t})}.
    \tag{F4}
\]
This is Zhao--Cui Eq. (23).  The correction is conceptually important: the TT
proposal can be biased, but if its conditional density is evaluable and covers
the target, importance weights can correct expectations.

\section{Backward Conditional Map, Path Estimation, And Smoothing}

The lower conditional map goes the other direction:
\[
    \widehat p_t(x_{t-1}\mid x_t,\alpha,y_{1:t})
    =
    \frac{\widehat q_t(x_t,\alpha,x_{t-1})}
    {\int \widehat q_t(x_t,\alpha,s)\dd s}.
    \tag{B1}
\]
The corresponding inverse conditional KR map samples
\[
    \widehat X_{t-1}^{(i)}
    =
    (F_{t,t-1}^{l})^{-1}
    \left(\Xi_{t-1}^{(i)}\mid \widehat X_t^{(i)},\widehat\alpha^{(i)}\right).
    \tag{B2}
\]
This is Zhao--Cui Eq. (24).

Starting at final time \(T\), draw
\[
    (\widehat X_T,\widehat\alpha,\widehat X_{T-1})
    \sim
    \widehat p_T(x_T,\alpha,x_{T-1}\mid y_{1:T}),
\]
then move backward using (B2).  The resulting approximate path density factors
as Zhao--Cui Eq. (25):
\[
    \widehat p(\alpha,x_{0:T}\mid y_{1:T})
    =
    \widehat p_T(x_T,\alpha,x_{T-1}\mid y_{1:T})
    \prod_{t=1}^{T-1}
    \widehat p_t(x_{t-1}\mid x_t,\alpha,y_{1:t}).
    \tag{B3}
\]
The exact smoothing posterior has the same Markov factorization form with exact
conditionals:
\[
    p(\alpha,x_{0:T}\mid y_{1:T})
    =
    p(x_T,\alpha,x_{T-1}\mid y_{1:T})
    \prod_{t=1}^{T-1}
    p(x_{t-1}\mid x_t,\alpha,y_{1:t}).
    \tag{B4}
\]
The approximation can therefore be corrected by an importance ratio.  The
unnormalized backward weight, Zhao--Cui Eq. (26), is
\[
    \omega_b(\alpha,x_{0:T})
    =
    \frac{
    p(\alpha)p_\alpha(x_0)
    \prod_{t=1}^{T}
    f_\alpha(x_t\mid x_{t-1})g_\alpha(y_t\mid x_t)}
    {
    \widehat p_T(x_T,\alpha,x_{T-1}\mid y_{1:T})
    \prod_{t=1}^{T-1}
    \widehat p_t(x_{t-1}\mid x_t,\alpha,y_{1:t})}.
    \tag{B5}
\]
The numerator is the exact unnormalized path posterior.  The denominator is the
proposal density generated by the backward TT maps.

\section{Error Propagation: What It Proves And What It Does Not Prove}

Section 4 of Zhao--Cui explains how local approximation errors propagate.  The
key identity is the triangle inequality.  Let \(\pi_t\) be the exact
unnormalized joint density of \((x_t,\alpha,x_{t-1})\) at time \(t\), \(q_t\)
the propagated density using the previous approximation, and
\(\widehat\pi_t\) the new TT approximation.  Then
\[
    D(\widehat\pi_t,\pi_t)
    \le
    D(\widehat\pi_t,q_t)+D(q_t,\pi_t).
    \tag{E1}
\]
The first term is the new fit error.  The second is inherited from the previous
time step.  Under bounded likelihood or bounded prediction-likelihood
conditions, the inherited error is multiplied by a finite constant.

For squared TT, Zhao--Cui use an \(L^2\) error on square roots because it
controls Hellinger distance after normalization.  The practical message is
precise but modest: if every local square-root approximation is accurate enough
and the model does not amplify errors too severely, the sequence of approximate
densities can remain controlled.  This is not a theorem that ranks will stay
small, that the companion implementation is globally differentiable, or that a
particular high-dimensional model will be accurate without diagnostics.

\section{Preconditioning}

The direct target \(q_t(x_t,\alpha,x_{t-1})\) can be sharply concentrated.  A TT
may then need high rank.  Zhao--Cui Section 5 introduces a change of variables
that makes the function closer to a product reference density before fitting.

Let
\[
    r=(x_t,\alpha,x_{t-1}),\qquad u=T_t(r),
\]
and let \(\eta(u)=\eta_1(u_1)\cdots\eta_D(u_D)\) be a product reference density.
Choose a bridging density \(\rho_t(r)\) that is easier to approximate than
\(q_t(r)\), and construct \(T_t\) so that the pushforward of \(\rho_t\) is close
to \(\eta\):
\[
    (T_t)_\#\rho_t(u)
    =
    \rho_t(T_t^{-1}(u))
    \left|\det\nabla T_t^{-1}(u)\right|
    \propto
    \eta(u).
    \tag{P1}
\]
This is Zhao--Cui Eq. (30).

The pushforward of \(q_t\) is
\[
    (T_t)_\#q_t(u)
    =
    q_t(T_t^{-1}(u))
    \left|\det\nabla T_t^{-1}(u)\right|.
\]
Using (P1), the Jacobian factor is proportional to
\(\eta(u)/\rho_t(T_t^{-1}(u))\), so the function to fit is
\[
    q^\sharp_t(u)
    =
    \frac{q_t(T_t^{-1}(u))}
    {\rho_t(T_t^{-1}(u))}
    \eta(u).
    \tag{P2}
\]
This is Zhao--Cui Eq. (31).  If \(\rho_t\) captures the main shape of \(q_t\),
then \(q_t/\rho_t\) is flatter than \(q_t\), and TT ranks may be smaller.

Fit a square-root TT in \(u\)-space:
\[
    \sqrt{q^\sharp_t(u)}\approx \phi^\sharp_t(u),
\]
and define
\[
    \widehat\nu^\sharp_t(u)
    =
    \phi^\sharp_t(u)^2+\tau^\sharp_t\eta(u).
    \tag{P3}
\]
This is Zhao--Cui Eq. (32).  Pulling the approximation back to physical
coordinates gives
\[
    \widehat\pi_t(r)
    =
    \frac{\widehat\nu^\sharp_t(T_t(r))}
    {\eta(T_t(r))}
    \rho_t(r).
    \tag{P4}
\]
This is Zhao--Cui Eq. (33).  The normalizer is unchanged by the change of
variables:
\[
    \int \widehat\pi_t(r)\dd r
    =
    \int \widehat\nu^\sharp_t(u)\dd u.
    \tag{P5}
\]

A simple bridge is Gaussian.  Estimate a mean and covariance for \(r\), set
\(\rho_t=\mathcal N(\mu_t,\Sigma_t)\), and use an affine Gaussian whitening map.
A nonlinear bridge uses powers of the three factors in \(q_t\):
\[
    \rho_t(r)
    =
    \widehat p_{t-1}(x_{t-1},\alpha)^{\beta_\pi}
    f_\alpha(x_t\mid x_{t-1})^{\beta_f}
    g_\alpha(y_t\mid x_t)^{\beta_g},
    \qquad
    \beta_\pi,\beta_f,\beta_g\in[0,1].
    \tag{P6}
\]
This is Zhao--Cui Eq. (34).  A squared-TT approximation of this bridge is
\[
    \widehat\rho_t(r)
    =
    \phi_{\rho,t}(r)^2+\tau_{\rho,t}\lambda(r).
    \tag{P7}
\]
This is Zhao--Cui Eq. (35).  The same algebra as (P4) applies with
\(\rho_t\) replaced by \(\widehat\rho_t\).

\section{Implementable Branch Objects And Data Structures}

A reader implementing a minimal method must store more than an equation.  At
time \(t\), a squared-TT filter object contains:
\[
    \mathcal B_t
    =
    \{
    \Omega_{t,k},\Psi_{t,k},\psi_{t,k,\ell},C_{t,k},
    R_{t,k},B_{t,k},\lambda_{t,k},\tau_t,c_t,
    \text{mass contractions},\text{map data}
    \}.
\]
Here \(\Omega_{t,k}\) is the coordinate domain, \(\Psi_{t,k}\) maps reference
coordinates to physical coordinates, \(\psi_{t,k,\ell}\) are basis functions,
\(C_{t,k}\) are TT cores, \(R_{t,k}\) are ranks, \(B_{t,k}\) are mass matrices,
\(\lambda_{t,k}\) are reference densities, \(\tau_t\) is the defensive mass, and
\(c_t\) is any log-scaling shift used during fitting.

For the adaptive Zhao--Cui algorithm these objects may change with the data,
with samples, and with approximation decisions.  For a differentiable
same-scalar target, the realized branch must be frozen before differentiation.

\subsection{What A Core Object Looks Like}

For coordinate \(k\), store
\[
    C_k\in\R^{R_{k-1}\times p_k\times R_k}.
\]
Given a scalar coordinate value \(r_k\), first compute the basis vector
\[
    b_k(r_k)=
    \bigl(\psi_{k,0}(r_k),\ldots,\psi_{k,p_k-1}(r_k)\bigr)^\top.
\]
Then form the matrix
\[
    H_k(r_k)[a,b]
    =
    \sum_{\ell=0}^{p_k-1}C_k[a,\ell,b]\,b_k(r_k)[\ell].
    \tag{D1}
\]
Evaluation of the TT is the loop
\[
    v_0=1,\qquad
    v_k=v_{k-1}H_k(r_k),\quad k=1,\ldots,D,
    \qquad
    \widehat h(r)=v_D.
    \tag{D2}
\]
The square-density object additionally stores mass contractions
\[
    M_{>k}
    =
    \int H_{k+1}(r_{k+1})\cdots H_D(r_D)
    H_D(r_D)^\top\cdots H_{k+1}(r_{k+1})^\top
    \dd r_{k+1:D}.
    \tag{D3}
\]
These matrices are enough to evaluate marginals from the left.  Analogous
left-side mass contractions evaluate reverse-order marginals.

\subsection{What Must Be Saved For The Next Time Step}

After time \(t\), the next step needs to evaluate
\[
    \widehat p_t(x_t;\beta)
    \quad\text{and}\quad
    \partial_{\beta_i}\widehat p_t(x_t;\beta)
\]
at future fitting points.  It is not enough to store a cloud of samples or only
\(\widehat Z_t\).  A fixed-branch derivative implementation stores:
\[
    \left\{
    C_{t,k},\dot C_{t,k}^{(i)},M_{t,>k},\dot M_{t,>k}^{(i)},
    \widehat Z_t,\dot{\widehat Z}_t^{(i)},
    \Psi_t,\lambda_t,\tau_t
    \right\}.
    \tag{D4}
\]
The derivative of the normalized marginal has the quotient form
\[
    \partial_{\beta_i}\widehat p_t(x_t;\beta)
    =
    \frac{\partial_{\beta_i}a_t(x_t;\beta)}{\widehat Z_t(\beta)}
    -
    \frac{a_t(x_t;\beta)\partial_{\beta_i}\widehat Z_t(\beta)}
    {\widehat Z_t(\beta)^2},
    \tag{D5}
\]
where
\[
    a_t(x_t;\beta)
    =
    \int \widehat q_t(x_t,x_{t-1};\beta)\dd x_{t-1}
    \tag{D6}
\]
is the unnormalized retained numerator.  Equation (D5) is the carried-filter
derivative needed in (G2) at the next time step.

\begin{lemma}[Carried filter object feeds the next target]
Fix the branch at time \(t\).  Suppose the stored object contains
\[
    \mathcal F_t
    =
    \{a_t,\dot a_t^{(i)},\widehat Z_t,\dot{\widehat Z}_t^{(i)},
      \Psi_t,\lambda_t,\tau_t\}.
\]
Define \(\widehat p_t\) and \(\dot{\widehat p}_t^{(i)}\) by (D5).  Then the
next transformed target and its derivative at any next-step fitting point are
computable from \(\mathcal F_t\), the model density evaluators, and their
derivatives.
\end{lemma}

\begin{proof}
At time \(t+1\), the fixed-branch target has the form
\[
    \widetilde q_{t+1}(z;\beta)
    =
    \widehat p_t(x_t(z);\beta)
    f_\beta(x_{t+1}(z)\mid x_t(z))
    g_\beta(y_{t+1}\mid x_{t+1}(z))
    J_{t+1}(z).
\]
The stored object evaluates \(\widehat p_t\).  The model supplies \(f_\beta\)
and \(g_\beta\), and the branch supplies \(J_{t+1}\).  Differentiating gives
the product-rule expression (G2) with \(t\) replaced by \(t+1\).  The first
term uses \(\dot{\widehat p}_t^{(i)}\), which is the quotient derivative in
(D5).  No old samples or external source-paper objects are needed.
\end{proof}

\section{One Concrete Branch, Fully Instantiated}

This section gives a single implementation lane.  It is intentionally narrower
than Zhao--Cui's adaptive code.  Its purpose is to remove ambiguity.  A coder
can implement this lane first, then replace pieces with adaptive or higher
performance choices later.

\subsection{Coordinate Order And Domain}

For the fixed-parameter likelihood case, use the two-block coordinate order
\[
    r_t=(x_t,x_{t-1})\in\R^2
\]
for a scalar state.  For vector states, concatenate coordinates as
\[
    r_t=(x_{t,1},\ldots,x_{t,m_x},x_{t-1,1},\ldots,x_{t-1,m_x}).
\]
Choose a fixed interval for each coordinate,
\[
    r_{t,k}\in[\ell_{t,k},u_{t,k}],
\]
and map from reference coordinate \(z_k\in[-1,1]\) by
\[
    r_{t,k}=a_{t,k}+h_{t,k}z_k,\qquad
    a_{t,k}=\frac{u_{t,k}+\ell_{t,k}}2,\qquad
    h_{t,k}=\frac{u_{t,k}-\ell_{t,k}}2.
    \tag{C1}
\]
The Jacobian is constant:
\[
    J_t=\prod_{k=1}^{D}h_{t,k}.
    \tag{C2}
\]
The branch stores \(\ell_{t,k},u_{t,k},a_{t,k},h_{t,k}\).  In the same-scalar
derivative these are constants and \(\dot J_t=0\).

\subsection{Basis, Points, Weights, Ranks}

Use normalized Legendre basis with the same degree \(p\) in every coordinate:
\[
    b_k(z_k)=
    \left(
    \sqrt{\frac12}P_0(z_k),
    \sqrt{\frac32}P_1(z_k),
    \ldots,
    \sqrt{\frac{2p-1}{2}}P_{p-1}(z_k)
    \right)^\top.
    \tag{C3}
\]
Use fixed Halton-type fitting points.  Let \(p_k^\star\) be the \(k\)-th prime.
If
\[
    j=\sum_{m=0}^{M}d_m(p_k^\star)^m,\qquad
    d_m\in\{0,\ldots,p_k^\star-1\},
\]
define
\[
    \operatorname{radinv}_{p_k^\star}(j)
    =
    \sum_{m=0}^{M}d_m(p_k^\star)^{-(m+1)},
    \qquad
    z_k^{(j)}=2\operatorname{radinv}_{p_k^\star}(j)-1.
    \tag{C4}
\]
Use \(j=1,\ldots,N_{\rm fit}\), weights
\[
    W=\frac1{N_{\rm fit}}I,
    \tag{C5}
\]
and fixed ranks
\[
    R_0=R_D=1,\qquad R_1,\ldots,R_{D-1}\ \text{declared before fitting}.
    \tag{C6}
\]
No pivot, point, rank, or domain choice is allowed to change during a
same-scalar derivative check.

\subsection{Defensive Reference, Mass, And Shift}

Use the uniform reference on \([-1,1]^D\):
\[
    \lambda(z)=2^{-D}.
    \tag{C7}
\]
Choose a defensive floor \(\epsilon_\tau>0\) and set
\[
    \tau_t=2^D\epsilon_\tau.
    \tag{C8}
\]
Then \(\tau_t\lambda(z)=\epsilon_\tau\).  Choose the stabilizing shift at the
base parameter \(\beta_0\):
\[
    c_t=-\max_{1\le j\le N_{\rm fit}}
    \log \widetilde q_t(z^{(j)};\beta_0).
    \tag{C9}
\]
The fitted square-root target is
\[
    y_j(\beta)=\exp(c_t/2)\sqrt{\widetilde q_t(z^{(j)};\beta)}.
    \tag{C10}
\]
For same-branch finite differences, \(c_t\) is reused for
\(\beta_0+h\) and \(\beta_0-h\).  Recomputing \(c_t\) changes the scalar.

\subsection{Boxed Convention: Shifted Fit, Unshifted Density}

The concrete branch uses this convention everywhere downstream.
\[
\boxed{
\begin{aligned}
&\text{Fitted TT:}\qquad
  \phi_t(z;\beta)\approx e^{c_t/2}\sqrt{\widetilde q_t(z;\beta)},\\
&\text{Approximate joint density:}\qquad
  \widehat q_t(z;\beta)
  =
  e^{-c_t}\phi_t(z;\beta)^2+\tau_t\lambda_t(z),\\
&\text{Retained numerator:}\qquad
  a_t(z_t;\beta)
  =
  \int
  \left[e^{-c_t}\phi_t(z_t,z_{t-1};\beta)^2
        +\tau_t\lambda_t(z_t,z_{t-1})\right]\dd z_{t-1},\\
&\text{Normalizer:}\qquad
  \widehat Z_t(\beta)
  =
  \int \widehat q_t(z;\beta)\dd z
  =
  e^{-c_t}\int\phi_t(z;\beta)^2\dd z+\tau_t.
\end{aligned}}
\tag{C11}
\]
Every carried filter, proposition, derivative, and finite-difference check uses
(C11).  If an implementation instead absorbs \(e^{-c_t/2}\) into the TT cores,
then it must remove the explicit \(e^{-c_t}\) factors everywhere.  Mixing the
two conventions changes the scalar.

\subsection{Initialization And Sweep Order}

Initialize every core to zero except the constant channel:
\[
    C_k[1,0,1]=1,\qquad k=2,\ldots,D,
\]
and
\[
    C_1[1,0,1]=\bar y,\qquad
    \bar y=\frac1{N_{\rm fit}}\sum_{j=1}^{N_{\rm fit}}y_j.
\]
If a rank is larger than one, all non-first rank channels start at zero.  This
plain initialization gives a constant initial approximation and avoids a random
initialization branch.

One bidirectional sweep means update cores \(1,2,\ldots,D\), then
\(D,D-1,\ldots,1\).  Fix the number of bidirectional sweeps \(S\).  The core
update is the ridge solve (A1-2d).  The branch stores the sweep order, \(S\),
\(\rho\), and the final cores.

\subsection{Pseudocode For One Forward Step}

\begin{enumerate}[leftmargin=2em]
\item Build reference fitting points \(z^{(j)}\) by (C4).
\item Map them to physical points \(r^{(j)}=\Psi_t(z^{(j)})\).
\item Evaluate
\[
    \widetilde q_{t,j}
    =
    \widehat p_{t-1}(x_{t-1}^{(j)};\beta)
    f_\beta(x_t^{(j)}\mid x_{t-1}^{(j)})
    g_\beta(y_t\mid x_t^{(j)})J_t.
\]
\item Form \(y_j=\exp(c_t/2)\sqrt{\widetilde q_{t,j}}\).
\item Initialize TT cores by the constant-channel rule.
\item For each fixed sweep and each core \(k\), form environments
\[
    L_{j,k}=H_1(z_1^{(j)})\cdots H_{k-1}(z_{k-1}^{(j)}),
\]
\[
    R_{j,k}=H_{k+1}(z_{k+1}^{(j)})\cdots H_D(z_D^{(j)}),
\]
then form \(A_{j,k}\) by (A1-2b) and solve (A1-2d).
\item After fitting, compute square-mass contractions by (M5)--(M6).
\item Compute
\[
    R_t=\int\phi_t(z)^2\dd z,\qquad
    \widehat Z_t=e^{-c_t}R_t+\tau_t.
\]
\item Build the carried numerator \(a_t\) by integrating old-state coordinates
using the same convention:
\[
    a_t(z_t;\beta)
    =
    \int
    \left[e^{-c_t}\phi_t(z_t,z_{t-1};\beta)^2
    +\tau_t\lambda_t(z_t,z_{t-1})\right]\dd z_{t-1}.
\]
\item Store \(a_t,\widehat Z_t\), branch objects, cores, and mass contractions.
\end{enumerate}

\subsection{Conditional-Map Inversion Protocol}

For a conditional coordinate \(r_k\), compute
\[
    F_k(s\mid r_{<k})
    =
    \int_{\ell_k}^{s}\widehat p(r_k'\mid r_{<k})\dd r_k'.
\]
To sample from the conditional, given \(u\in(0,1)\), solve
\[
    F_k(s\mid r_{<k})-u=0,\qquad s\in[\ell_k,u_k].
\]
Use bisection as the baseline because it only needs monotonicity.  Fix a
tolerance \(\varepsilon_{\rm root}\) and maximum iterations \(N_{\rm root}\).
The branch stores both.  The bisection update is:
\[
    m=\frac{a+b}{2},\qquad
    \text{if }F_k(m\mid r_{<k})<u\text{ set }a=m,\text{ else set }b=m.
\]
Stop when \(b-a\le\varepsilon_{\rm root}\) or the iteration count reaches
\(N_{\rm root}\).  Raise a failure diagnostic if the final residual
\(|F_k(m\mid r_{<k})-u|\) exceeds the declared tolerance.

\section{Fixed-Branch TT Filtering Recursion}

The fixed-branch BayesFilter variant chooses all discrete approximation objects
before differentiating.  It is narrower than the full Zhao--Cui adaptive code,
but it is the variant for which an analytical derivative can be stated without
changing the scalar under the derivative.

There is one important notational separation.  In the Zhao--Cui learning
algorithm, the unknown parameter \(\alpha\) may be carried as a random variable
inside the posterior density.  In an HMC-style likelihood calculation, the
parameter is instead an external argument that is varied by the sampler.  To
avoid using the same symbol for both roles, this section writes the external
differentiated parameter as
\[
    \beta\in\R^{m_\beta}.
\]
The fixed-branch likelihood variant conditions on a trial \(\beta\) and filters
only the latent state.  Its unnormalized update is
\[
    q_t(x_t,x_{t-1};\beta)
    =
    \widehat p_{t-1}(x_{t-1};\beta)
    f_\beta(x_t\mid x_{t-1})
    g_\beta(y_t\mid x_t).
\]
If one also wants to carry a learned static parameter as a random coordinate,
that coordinate can be appended to the state; the derivative below is still
with respect to the external \(\beta\), not with respect to an integration
variable.

Let \(z\in[-1,1]^D\) be reference coordinates and
\[
    r=\Psi_t(z)=(x_t,x_{t-1}).
\]
Let \(J_t(z)=|\det\nabla\Psi_t(z)|\).  The transformed target is
\[
    \widetilde q_t(z;\beta)
    =
    \widehat p_{t-1}(x_{t-1};\beta)
    f_\beta(x_t\mid x_{t-1})
    g_\beta(y_t\mid x_t)
    J_t(z).
    \tag{FB1}
\]
The branch fixes points \(z^{(j)}\), weights \(w_j\), basis functions, ranks,
and a core-construction algorithm.  The fitted square-root TT is
\[
    \phi_t(z;\beta)=G_{t,1}(z_1;\beta)\cdots G_{t,D}(z_D;\beta).
    \tag{FB2}
\]
The fixed-branch approximate density is
\[
    \widehat q_t(z;\beta)
    =
    e^{-c_t}\phi_t(z;\beta)^2+\tau_t\lambda_t(z),
    \qquad
    \int\lambda_t(z)\dd z=1.
    \tag{FB3}
\]
Its normalizer is
\[
    \widehat Z_t(\beta)
    =
    \int\widehat q_t(z;\beta)\dd z
    =
    e^{-c_t}\int\phi_t(z;\beta)^2\dd z+\tau_t.
    \tag{FB4}
\]
This is exactly the boxed convention (C11).
The next filter object is the normalized marginal
\[
    \widehat p_t(x_t;\beta)
    =
    \frac{
    \int \widehat q_t(z_t,z_{t-1};\beta)\dd z_{t-1}}
    {\widehat Z_t(\beta)}
    \left|\det\nabla z_t/\nabla x_t\right|.
    \tag{FB5}
\]
The Jacobian factor appears if the filter is evaluated in physical coordinates.
If the whole recursion is implemented in reference coordinates, it is absorbed
into the stored domain maps.

\begin{proposition}[Fixed-branch recursion is a normalized approximate filter]
Fix a time horizon \(T\), model densities \(f_\beta,g_\beta\), initial density
\(p_\beta(x_0)\), and fixed branch objects for each time step.  Suppose every
\(\widehat q_t(\cdot;\beta)\) defined by (FB3) is nonnegative, integrable, and
has \(0<\widehat Z_t(\beta)<\infty\).  Define \(\widehat p_t\) by (FB5),
starting from the fixed initial approximation.  Then each \(\widehat p_t\) is a
normalized density on the retained state variable \(x_t\).  Moreover, the
recursion is exactly the Bayes filtering recursion applied to the declared
approximate joint density \(\widehat q_t\), not to the exact joint density.
\end{proposition}

\begin{proof}
Nonnegativity follows from
\(\phi_t^2\ge0\), \(\tau_t>0\), and \(\lambda_t\ge0\).  Since
\(\widehat Z_t=\int\widehat q_t\) is positive and finite, the normalized joint
density
\[
    \widehat p_t^{\rm joint}(z)
    =
    \widehat q_t(z)/\widehat Z_t
\]
integrates to one.  The retained filter is the marginal of this normalized
joint density:
\[
    \int \widehat p_t(z_t)\dd z_t
    =
    \int
    \left[
    \int \widehat p_t^{\rm joint}(z_t,z_{t-1})\dd z_{t-1}
    \right]\dd z_t
    =
    \int \widehat p_t^{\rm joint}(z)\dd z
    =
    1.
\]
The change-of-variables factor in physical coordinates is the standard
Jacobian factor for density transformation, so normalization is preserved.  The
construction uses the exact Bayes filtering algebra with \(q_t\) replaced by
the declared approximation \(\widehat q_t\).  No step asserts that
\(\widehat q_t=q_t\), so the result is a normalized approximate filtering
recursion rather than an exact-posterior theorem.
\end{proof}

\section{Same-Scalar Analytical Derivative}

For parameter inference one may want the approximate log evidence
\[
    \widehat\ell_T(\beta)
    =
    \sum_{t=1}^T \log \widehat Z_t(\beta).
    \tag{G1}
\]
The word ``same'' matters.  The derivative must be the derivative of the scalar
actually computed by the forward recursion: same domains, same points, same
ranks, same defensive mass, same scaling shifts, same linear solves, and same
stored filter objects.

Let a dot denote \(\partial/\partial\beta_i\).  At fitting point \(z^{(j)}\),
the transformed target derivative is
\[
\begin{aligned}
 \dot{\widetilde q}_{t,j}
 &=
 J_{t,j}
 \Big[
 \dot{\widehat p}_{t-1}(x_{t-1}^{(j)};\beta)
 f_{\beta,j}g_{\beta,j}
 + \widehat p_{t-1,j}\dot f_{\beta,j}g_{\beta,j}
 + \widehat p_{t-1,j}f_{\beta,j}\dot g_{\beta,j}
 \Big]
 + \widehat p_{t-1,j}f_{\beta,j}g_{\beta,j}\dot J_{t,j}.
\end{aligned}
\tag{G2}
\]
For a frozen domain map, \(\dot J_{t,j}=0\).  If the fitted target is
\[
    y_{t,j}=e^{c_t/2}\sqrt{\widetilde q_{t,j}},
\]
with frozen \(c_t\), then
\[
    \dot y_{t,j}
    =
    \frac{e^{c_t/2}}{2\sqrt{\widetilde q_{t,j}}}
    \dot{\widetilde q}_{t,j},
    \tag{G3}
\]
where a positive floor is needed in implementation if \(\widetilde q_{t,j}\) is
near zero.

For a fixed least-squares core update, write the normal equations as
\[
    N_k g_k=d_k,
    \qquad
    N_k=A_k^\top W A_k+\rho I,
    \qquad
    d_k=A_k^\top W y.
    \tag{G4}
\]
Here \(g_k=\vecop(C_k)\).  Differentiating gives
\[
    N_k\dot g_k
    =
    \dot d_k-\dot N_k g_k,
    \tag{G5}
\]
with
\[
    \dot N_k
    =
    \dot A_k^\top W A_k+A_k^\top W\dot A_k,
    \qquad
    \dot d_k
    =
    \dot A_k^\top Wy+A_k^\top W\dot y.
    \tag{G6}
\]
The design matrix derivative \(\dot A_k\) comes from the left and right TT
environments, which depend on already differentiated neighboring cores during
the fixed sweep order.  Thus the derivative follows the same ordered sequence
of stored linear solves as the forward core construction.

\subsection{Differentiated Environments In Indices}

For one fitting point \(z^{(j)}\), define
\[
    L_{j,k}=H_1(z_1^{(j)})\cdots H_{k-1}(z_{k-1}^{(j)}),
    \qquad
    R_{j,k}=H_{k+1}(z_{k+1}^{(j)})\cdots H_D(z_D^{(j)}).
\]
The derivative of the left environment is propagated by
\[
    \dot L_{j,1}=0,
\]
\[
    \dot L_{j,k+1}
    =
    \dot L_{j,k}H_k(z_k^{(j)})
    +
    L_{j,k}\dot H_k(z_k^{(j)}).
    \tag{G6a}
\]
The derivative of the right environment is propagated by
\[
    \dot R_{j,D}=0,
\]
\[
    \dot R_{j,k-1}
    =
    \dot H_k(z_k^{(j)})R_{j,k}
    +
    H_k(z_k^{(j)})\dot R_{j,k}.
    \tag{G6b}
\]
If the basis and points are frozen, then
\[
    \dot H_k(z_k^{(j)})[a,b]
    =
    \sum_{\ell=0}^{p_k-1}
    \dot C_k[a,\ell,b]\,b_k(z_k^{(j)})[\ell].
    \tag{G6c}
\]
Thus the design-row derivative is not mysterious.  Since
\[
    A_{j,k}=R_{j,k}^\top\otimes b_k(z_k^{(j)})^\top\otimes L_{j,k},
\]
and \(b_k\) is fixed on the branch,
\[
    \dot A_{j,k}
    =
    \dot R_{j,k}^{\top}\otimes b_k(z_k^{(j)})^\top\otimes L_{j,k}
    +
    R_{j,k}^{\top}\otimes b_k(z_k^{(j)})^\top\otimes \dot L_{j,k}.
    \tag{G6d}
\]
Equations (G6a)--(G6d) are the operational recipe for forming
\(\dot A_k\) in (G6).

The normalizer derivative follows from the mass contraction.  If
\[
    R_t(\beta)=\int \phi_t(z;\beta)^2\dd z,
    \qquad
    \widehat Z_t(\beta)=e^{-c_t}R_t(\beta)+\tau_t,
\]
and \(c_t,\tau_t\) are frozen, then
\[
    \dot{\widehat Z}_t=e^{-c_t}\dot R_t
    =
    2e^{-c_t}\int \phi_t(z;\beta)\dot\phi_t(z;\beta)\dd z.
    \tag{G7}
\]
In core form this is computed by differentiating the same square-core mass
contractions used to compute \(R_t\).

\subsection{Differentiated Mass Contraction In Indices}

The right mass recursion was
\[
    M_{>j-1}[a,a']
    =
    \sum_{b,b',\ell,\ell'}
    C_j[a,\ell,b]M_{>j}[b,b']C_j[a',\ell',b']B_j[\ell,\ell'].
\]
On a frozen branch \(B_j\) is fixed.  Differentiating gives
\[
\begin{aligned}
 \dot M_{>j-1}[a,a']
 &=
 \sum_{b,b',\ell,\ell'}
 \dot C_j[a,\ell,b]M_{>j}[b,b']C_j[a',\ell',b']B_j[\ell,\ell']\\
 &\quad+
 \sum_{b,b',\ell,\ell'}
 C_j[a,\ell,b]\dot M_{>j}[b,b']C_j[a',\ell',b']B_j[\ell,\ell']\\
 &\quad+
 \sum_{b,b',\ell,\ell'}
 C_j[a,\ell,b]M_{>j}[b,b']\dot C_j[a',\ell',b']B_j[\ell,\ell'].
\end{aligned}
\tag{G7a}
\]
Starting from \(\dot M_{>D}=0\), this recursion produces the derivative of the
complete square contraction.  With the shifted-target convention,
\[
    \widehat Z_t=e^{-c_t}R_t+\tau_t,
    \qquad
    \dot{\widehat Z}_t=e^{-c_t}\dot R_t,
    \tag{G7b}
\]
because \(c_t\) and \(\tau_t\) are frozen branch constants.

Therefore
\[
    \partial_i\widehat\ell_T(\beta)
    =
    \sum_{t=1}^T
    \frac{\partial_i\widehat Z_t(\beta)}
    {\widehat Z_t(\beta)}.
    \tag{G8}
\]

\begin{proposition}[Fixed-branch gradient differentiates the declared scalar]
Assume all branch objects are fixed; all model densities, basis evaluations, and
domain maps used on the branch are differentiable in \(\beta\); every fitted
core is obtained by a fixed finite sequence of ridge-regularized linear solves
with nonsingular matrices \(N_k\); and
\(0<\widehat Z_t(\beta)<\infty\).  Then the recursion (G2)--(G8) computes
\(\nabla_\beta\widehat\ell_T(\beta)\) for the scalar in (G1), where
\(\widehat Z_t\) is the normalizer computed by the fixed-branch forward pass.
\end{proposition}

\begin{proof}
The proof is induction over time and over the fixed core-update sequence.  At
time \(0\), the initial filter evaluator and its derivative are assumed fixed
and differentiable.  Suppose the previous filter evaluator and derivative are
available.  Then the transformed target derivative (G2) follows from the
product rule and the chain rule.  The square-root target derivative (G3)
follows from differentiating \(e^{c_t/2}\sqrt{\widetilde q}\), with
\(\dot c_t=0\) on the fixed branch.

For each core update, the forward pass solves \(N_k g_k=d_k\).  Because
\(N_k\) is nonsingular and differentiable, differentiating
\(N_k g_k=d_k\) yields
\[
    \dot N_k g_k+N_k\dot g_k=\dot d_k,
\]
and hence (G5).  The fixed sweep order composes finitely many differentiable
maps, so all fitted cores are differentiable functions of \(\beta\) on the
branch.

The square contraction \(R_t=\int\phi_t^2\) is a finite contraction of
differentiable core coefficients and fixed mass matrices.  Differentiating that
contraction and then applying the frozen scaling convention
\(\widehat Z_t=e^{-c_t}R_t+\tau_t\) gives (G7).  Since
\(\dot c_t=0\) and \(\dot\tau_t=0\), no additional terms appear.  Finally,
differentiating
\(\widehat\ell_T=\sum_t\log\widehat Z_t\) gives (G8).  Each term is computed
from the same stored objects used by the forward pass, so the derivative is of
the declared scalar, not of a nearby adaptive algorithm.
\end{proof}

\section{Numerical Stabilization, Failure Diagnostics, And Finite-Difference Test}

A minimal implementation should record the following diagnostics at every time:
\[
    \widehat Z_t,\qquad
    \frac{\tau_t}{\widehat Z_t},\qquad
    \max_k R_{t,k},\qquad
    \kappa(N_k),\qquad
    \frac{\|A g-y\|_W}{\|y\|_W},\qquad
    \min_j \widetilde q_{t,j}.
\]
The defensive mass ratio \(\tau_t/\widehat Z_t\) detects whether the floor is
dominating.  The rank and condition numbers detect whether the TT fit is too
hard for the chosen branch.  The residual detects local fitting failure.  The
minimum target value detects square-root instability.

For the gradient, use same-branch finite differences.  Choose a base
\(\beta_0\), build and freeze the branch, then evaluate
\[
    \frac{\widehat\ell_T(\beta_0+h e_i;B)
    -
    \widehat\ell_T(\beta_0-h e_i;B)}
    {2h}
\]
with exactly the same branch \(B\).  Compare it to the analytical derivative
from (G8).  Rebuilding the branch at \(\beta_0\pm h e_i\) tests a different
adaptive scalar and is not a valid same-scalar derivative test.

\section{Minimal Runnable Example Specification}

A minimal example should use a one-dimensional nonlinear state model with one
unknown scalar parameter:
\[
    x_t=\beta x_{t-1}+\sigma\epsilon_t,\qquad
    y_t=\sin(x_t)+\eta_t,
\]
with Gaussian noises.  Use \(T=2\) so the second step must evaluate the carried
filter from step 1.  Use a fixed reference interval, fixed Legendre basis,
fixed fitting points, fixed rank one or two, fixed ridge, fixed \(\tau_t\), and
fixed scaling shifts.  The script should print:
\[
    \widehat\ell_T(\beta_0),\qquad
    \partial_\beta\widehat\ell_T(\beta_0),\qquad
    \text{central finite-difference error}.
\]
Passing this example does not prove high-dimensional performance.  It proves
that the value path and derivative path are aligned for the declared scalar.

\subsection{Two-Step Value Path}

The two-step example should be written so the second step cannot cheat.  At
time \(1\), build
\[
    q_1(x_1,x_0;\beta)
    =
    p_\beta(x_0)f_\beta(x_1\mid x_0)g_\beta(y_1\mid x_1).
\]
Fit \(\phi_1\), form
\[
    \widehat q_1(z_1,z_0;\beta)
    =
    e^{-c_1}\phi_1(z_1,z_0;\beta)^2+\tau_1\lambda_1(z_1,z_0),
\]
and compute
\[
    \widehat Z_1(\beta)
    =
    e^{-c_1}\int \phi_1(z_1,z_0;\beta)^2\dd z_1\dd z_0+\tau_1.
\]
Then store the carried filter
\[
    \widehat p_1(x_1;\beta)
    =
    \frac{\int \widehat q_1(z_1,z_0;\beta)\dd z_0}
    {\widehat Z_1(\beta)}
    \left|\frac{\dd z_1}{\dd x_1}\right|.
\]
At time \(2\), the target must call this stored evaluator:
\[
    q_2(x_2,x_1;\beta)
    =
    \widehat p_1(x_1;\beta)
    f_\beta(x_2\mid x_1)g_\beta(y_2\mid x_2).
\]
The total scalar is
\[
    \widehat\ell_2(\beta)
    =
    \log\widehat Z_1(\beta)+\log\widehat Z_2(\beta).
\]
Here \(\widehat Z_2\) is computed by the same convention after fitting
\(\phi_2\) to the shifted square-root target for \(q_2\):
\[
    \widehat q_2(z_2,z_1;\beta)
    =
    e^{-c_2}\phi_2(z_2,z_1;\beta)^2+\tau_2\lambda_2(z_2,z_1),
    \qquad
    \widehat Z_2(\beta)
    =
    e^{-c_2}\int \phi_2(z_2,z_1;\beta)^2\dd z_2\dd z_1+\tau_2 .
\]
The derivative path must carry
\(\partial_\beta\widehat p_1(x_1;\beta)\) into the derivative of \(q_2\).
If a script recomputes step 1 differently when evaluating step 2, it is no
longer testing the fixed-branch recursion.

\section{Reader-Facing Conclusion}

Zhao--Cui's method is persuasive because it attacks the filtering recursion as
a sequence of function-approximation problems.  The posterior update creates a
three-block unnormalized density in \((x_t,\alpha,x_{t-1})\).  A tensor train
tries to store that function without a full grid.  Squaring a TT for the
square-root target makes the approximation nonnegative.  Mass contractions give
normalizers and marginals.  Marginal ratios give conditional densities.  Those
conditional densities give KR maps for forward filtering proposals and backward
smoothing proposals.

For BayesFilter's gradient-driven use case, the adaptive research algorithm
must be separated from a fixed-branch scalar.  Once the branch is fixed, the
approximate normalizers define a clear scalar
\(\widehat\ell_T=\sum_t\log\widehat Z_t\), and ordinary calculus differentiates
that scalar through target evaluations, TT core solves, contractions, and
carried filters.  The result is not a proof of exact posterior accuracy.  It is
a mathematically honest route from Zhao--Cui's sequential TT idea to an
implementable approximate likelihood and same-scalar derivative.

\appendix
\section{Equation And Algorithm Disposition Summary}

The detailed ledger beside this note gives every numbered equation and
algorithm disposition.  In the main text, Zhao--Cui Eqs. (1)--(26) and
(30)--(35), and Algorithms 1--5, are either reconstructed directly or discussed
as support for the reconstruction.  Numerical example equations from Section 6
are used only as evidence of demonstrated examples, not as theorem support.

\end{document}
